% Progress report -- Skill Progress Monitoring
% Build:  pdflatex progress_report && bibtex progress_report && pdflatex progress_report && pdflatex progress_report
% Target: 2 pages. See docs/report/README.md before adding anything.

\documentclass[11pt,a4paper]{article}

\usepackage[margin=2.0cm]{geometry}
\usepackage{lmodern}          % scalable fonts: microtype's expansion needs them
\usepackage[T1]{fontenc}
\usepackage[utf8]{inputenc}
\usepackage{enumitem}
\usepackage[compact,small]{titlesec}
\usepackage{microtype}
\usepackage{amsmath}
\usepackage[hidelinks]{hyperref}

\setlist{nosep,leftmargin=*}
\titlespacing*{\section}{0pt}{0.9ex plus .1ex}{0.3ex}
\newcommand{\TODO}[1]{\textbf{[TODO: #1]}}
\let\oldthebibliography\thebibliography
\renewcommand{\thebibliography}[1]{\oldthebibliography{#1}%
  \setlength{\itemsep}{0pt}\setlength{\parsep}{0pt}}

\title{\vspace{-1.2cm}Skill Progress Monitoring\\[2pt]\large Progress Report}
\author{Pinhas Nehoray Aburmad\\
  \small Supervisor: Prof.\ Ronen Brafman\\
  \small Department of Computer Science, Ben-Gurion University of the Negev}
\date{\small 7 September 2026}

\begin{document}
\maketitle
\vspace{-1.0cm}

\section{The claim}

Monitoring whether a robot skill is progressing, and where it is going wrong, normally
requires a specification written by hand for each skill and each robot. The claim of
this work is that such a monitor can instead be \emph{synthesised} from a free-language
description of the skill, with the engine unchanged across skills and across
embodiments.

The obstacle is grounding. A language model asked to write a temporal-logic
specification from a description of a skill has no way of knowing which quantities the
robot can actually observe, so it writes conditions over terms that do not exist on the
target platform. An ungrounded specification is not a weak monitor; it is not a monitor
at all. Existing work addresses this by grounding into a signature authored by hand and
fixed before inference, or by judging groundedness with a learned
classifier~\cite{english2025ginsign}.

The proposal here is that grounding can target the robot's own declared sensor schema,
live, and be checked mechanically. Each embodiment publishes an \emph{adapter} saying
what it can sense; each skill has a \emph{specification}; and a validation oracle rejects
any generated specification referring to anything outside that schema. Because this is a
free-variable check rather than a judgement, it cannot admit a specification that would
fail on the robot.

Three things would settle the claim: that one engine runs unchanged across skills and
embodiments; that specifications generated this way execute on the robot they were
generated for; and that the resulting monitor detects real failures, with intervention on
its verdicts improving outcomes.

\section{How it works}

Progress is tracked by a phase machine over temporal-logic milestones, while named
failure modes run as parallel safety monitors, so a fault can be attributed to a
condition with a name rather than to the specification as a whole. Automata are built
with Spot~3~\cite{duretlutz2022spot} and verdicts follow the three-valued semantics of
Bauer et al.~\cite{bauer2011runtime}. Responses are graded rather than binary, from a
warning through to an abort. A compiled specification for the navigation skill reads, in
abbreviated form:

\begin{center}\footnotesize
\begin{tabular}{@{}ll@{}}
progress & \texttt{F(mission\_started \& F(path\_active \& F(moving\_to\_target \& F(mission\_finished))))} \\
safety & \texttt{G(!collision\_risk)}\qquad \texttt{G(upright)} \\
\end{tabular}
\end{center}

\noindent
Each proposition carries an executable rule over the adapter's declared sensor keys. The
implementation is at \url{https://github.com/pinhas019/LtlMonitor}.

\section{Progress this period}

\textbf{MiniGrid and MAAOS.} The continuous-physics layer over MiniGrid carries the
multi-agent cooperative box-push environment the monitor runs against. Three
specifications are loaded: one task-level, and one for each MAAOS skill. The monitor
watches a coordination sequence (reach the box, align on one side, push together, place
it on target), safety guards around lava and loss of the push formation, and a deadlock
condition on both agents' poses. Adding a skill is a specification file and nothing else;
adding a hazard type still requires a change in the environment as well.

\textbf{The robot.} The system was deployed across two tiers --- clock, evaluator and
recorder on the Unitree G1, monitor and operator console on the development machine ---
and run on the robot. The recorded episode replayed on the development machine to
identical verdicts throughout, which is the test for the claim that the monitor does not
depend on the machine it runs on. That is the first result on hardware, and a limited
one: the run exercised only part of the specification.

\textbf{Navigation application.} I am working with Elias on the navigation application
intended for deployment on the G1.

\textbf{High-fidelity simulation.} Under development. Sensor bridging for Isaac is in
place. The purpose is to re-execute episodes as counterfactuals rather than to visualise
them, so the work is gated on fidelity validation against recorded runs. What remains is
the robot's own navigation stack inside the simulator, and scenario reset.

\textbf{Literature.} Ongoing. Twenty papers have been read closely and summarised so
far. The review has settled how this work is positioned against the nearest prior art and
produced a ranked list of corrections to our own implementation.

\section{How far the claim has got}

Against the first condition, one engine now runs unchanged across three embodiments ---
the real G1, MuJoCo and Isaac Lab expose an identical sensor schema over entirely
different topics --- and across more than one skill, since adding a skill is a
specification file alone.

The third condition has its first evidence, from the cooperative box-push work. On an
injected-deadlock scenario the unmonitored run exhausted its whole 80-step budget and
failed in every seed; detection alone halted it at step 10, and halting followed by a
single re-plan turned the same scenario into a success in every seed, in 45 steps, with
no false alarms on clean runs or runs with non-fatal hazards. A second case bears on the
second condition: running the language-model-evaluated variant against the real push
skill showed that two predicates demanded a formation that skill never sustains. The
monitor was flagging correct behaviour and the fault lay in the generated specification
--- exactly what the oracle exists to catch before a specification runs at all. These
runs drove scripted reference plans rather than the live MAAOS planner, so they establish
detection and the value of intervention, but not yet either under the planner.

Once the navigation application is deployed on the G1, the monitor will be used to assess
it: the first skill developed independently of the monitor, and so the sharpest available
test of the first two conditions.

\section{Next steps}

The pace is set by two things becoming available rather than by a calendar. The first is
the high-fidelity simulation: once it is ready, episodes can be driven automatically and
in volume, which is what the claim needs in order to be evaluated rather than
demonstrated. The second is a wider set of skills --- which is where the navigation
application and further MAAOS skills come in.

Alongside those, we intend to compare this monitor against existing runtime-verification
tools for ROS~\cite{ferrando2020rosmonitoring,rtamt2024}, on the axes where we expect to
differ: synthesis of the specification from a free-language description, grounding in a
declared sensor schema, and portability of one specification across embodiments.

\bibliographystyle{abbrv}
{\footnotesize\bibliography{references}}

\end{document}
